\Qhead{2.8}{Волны в нелинейных средах без дисперсии}{Волны в нелинейных средах без дисперсии. Образование разрывов. Ударные волны. Уравнение Бюргерса для диссипативной среды и свойства его решений. Генерация гармоник исходного монохроматического сигнала, эффекты нелинейного поглощения, насыщения и детектирования.}
\Subfile{2.8.1}{sections/02_waves/q08/s01}
\Subfile{2.8.2}{sections/02_waves/q08/s02}

\recap{Амплитудно-зависимая скорость искривляет профиль и приводит к градиентной катастрофе. Диссипация формирует конечную толщину фронта. Гармоники переносят энергию в область сильного поглощения; разностные частоты дают детектирование.}
\board{Характеристики простой волны, расстояние разрыва, Бюргерс, tanh-фронт и энергетический баланс.}
\oral{Почему консервативная нелинейность может усиливать наблюдаемое поглощение в диссипативной среде?}
\bait{Можно оставить механизм насыщения: крутой фронт создаёт высокие гармоники, которые быстрее поглощаются.}
\sources{Цысарь~\cite{tsysar2}, лекции~1--3, PDF-с.~5--17; Цысарь~\cite{tsysar1}, лекции~3--4. Бюргерс и преобразование Коула–Хопфа приведены в явно нормированной форме; решение проверяется подстановкой.}
